%%=============================================================================
%% ANEXO B
%%=============================================================================
\chapter{Anexo: Derivaciones y tablas auxiliares}
\label{anexo:derivaciones}

\section{Sesgo del bit 7 en la extracción de \texorpdfstring{$n_p$}{np} módulo 256}

Sea $n_p \sim \text{Geom}(q)$ en $\{1, 2, \ldots\}$, con $\Pr[n_p = j] = q\, r^{j-1}$ y $r = 1 - q$. La probabilidad de que $n_p$ caiga en un intervalo $[a, b]$ es
\begin{equation}
    \Pr[a \leq n_p \leq b] = \sum_{j=a}^{b} q\, r^{j-1} = r^{a-1} - r^{b}.
\end{equation}
El bit 7 de $n_p \bmod 256$ vale 1 cuando $n_p \in [256\ell + 128,\; 256\ell + 255]$ para algún $\ell \geq 0$. Sumando sobre $\ell$,
\begin{equation}
    p_7 = \sum_{\ell = 0}^{\infty} \left( r^{256\ell + 127} - r^{256\ell + 255} \right) = \frac{r^{127} - r^{255}}{1 - r^{256}},
\end{equation}
que es la Ecuación~\eqref{eq:sesgo_bit7}. Para el XOR de $m$ valores independientes con la misma probabilidad $p_7$, el lema de apilamiento da
\begin{equation}
    \Pr[\text{XOR} = 1] = \frac{1}{2} - \frac{1}{2} \prod_{i=1}^{m} (1 - 2p_7) = \frac{1}{2} - \frac{1}{2}(1 - 2p_7)^m,
\end{equation}
que es la Ecuación~\eqref{eq:sesgo_xor}. Con $q = 1/120$ ($N = 5$), $p_7 = 0.25733$ y, con $m = 5$, $\Pr[\text{XOR} = 1] = 0.486536$. Con $q = 1/720$ ($N = 6$), $p_7 = 0.45628$ y, con $m = 4$, $\Pr[\text{XOR} = 1] = 0.499971$.

\section{Umbrales de las pruebas de salud}

La Tabla~\ref{tab:umbrales_salud} lista los umbrales que usa \texttt{qpp-rng-iot} según el presupuesto operativo conservador de min-entropía $h$, con $\alpha = 2^{-20}$ y una ventana de $W = 512$ muestras para la APT. Los valores se calcularon de forma exacta con la distribución binomial. Un valor de $h$ no tabulado se redondea hacia abajo al valor más cercano de la tabla, y uno mayor que 4 se trata como 4. Declarar menos entropía de la que realmente contiene la fuente eleva los umbrales, de modo que las pruebas toleran más repeticiones y resultan menos sensibles que las calculadas para el valor real; un valor de 0 se trata como medio bit. El parámetro $h$ constituye un presupuesto operativo adoptado a partir de estimaciones experimentales y no equivale por sí solo a una acreditación formal de la fuente.

\begin{table}[H]
\centering
\caption{Umbrales de las pruebas de salud según el presupuesto operativo de min-entropía}
\label{tab:umbrales_salud}
\begin{tabular}{@{}ccc@{}}
\toprule
\textbf{$h$ operativo (bits por muestra)} & \textbf{$C_{\text{RCT}}$} & \textbf{$C_{\text{APT}}$ ($W = 512$)} \\
\midrule
0 (tratado como 1/2) & 41 & 410 \\
1 & 21 & 311 \\
2 & 11 & 177 \\
3 & 8 & 103 \\
4 (máximo tabulado) & 6 & 62 \\
\bottomrule
\end{tabular}
\end{table}

\section{Muestras por byte}

La Tabla~\ref{tab:muestras_por_byte} lista las muestras que la variante IoT absorbe por byte, con 4 caminatas por byte, según el presupuesto operativo de min-entropía $h$. Los valores corresponden a la entropía presupuestada a la entrada de la tubería, condicionada a la validez de $h$, y no constituyen una garantía de min-entropía en los bytes de salida.

\begin{table}[H]
\centering
\caption{Muestras absorbidas por byte de salida}
\label{tab:muestras_por_byte}
\begin{tabularx}{\textwidth}{@{}>{\centering\arraybackslash}X>{\centering\arraybackslash}X>{\centering\arraybackslash}X>{\centering\arraybackslash}X@{}}
\toprule
\textbf{$h$ operativo} & \textbf{Muestras por caminata} & \textbf{Muestras por byte} & \textbf{Bits presupuestados por byte a la entrada} \\
\midrule
1 & 2 & 8 & 8 \\
2 & 1 & 4 & 8 \\
3 & 1 & 4 & 12 \\
4 & 1 & 4 & 16 \\
\bottomrule
\end{tabularx}
\end{table}